% !TEX root = Frozen Chicken Stock.tex
\documentclass[14pt]{extarticle}
\usepackage[utf8]{inputenc}
% --- REMOVE PAGE NUMBERING ---
\pagestyle{empty}
% --- MARGIN SETUP ---
\usepackage[left=0.75in, right=0.75in, top=0.5in, bottom=0.5in]{geometry}
% --- LIST SETTINGS ---
\usepackage{enumitem}
\setlist[itemize]{leftmargin=1em}
% Labels are centered in a fixed-width box so plain bullets line up with the
% wider prep icons (\veg, \measure) below.
\SetLabelAlign{center}{\makebox[\labelwidth][c]{#1}}
\setlist[itemize]{leftmargin=1em, itemsep=0.2em, label={\raisebox{-0.1ex}{\large\textbullet}},
    align=center, labelwidth=0.9em, labelsep=0.4em}
\setlist[enumerate]{leftmargin=1.5em, label=\textbf{\arabic*.}, itemsep=0.25em}
\usepackage{parskip}
% --- CONTROL SECTION FONT SIZE ---
\usepackage{titlesec}
\titleformat{\section}{\large\bfseries}{}{0em}{}
\titleformat{\subsection}{\normalsize\bfseries}{}{0em}{}
\titlespacing*{\subsection}{0pt}{0.75em}{0.25em}
\titlespacing*{\section}{0pt}{1em}{0.5em}
% --- FONT SETUP ---
\usepackage{helvet}
\renewcommand{\familydefault}{\sfdefault}
% --- CUSTOM COMMAND FOR TITLE ---
\newcommand{\makeRecipeTitle}[1]{
    {\noindent \Large \textbf{#1}}
    \par \vspace{-0.75em}
}
% --- CUSTOM COMMAND FOR SERVINGS ---
% Usage: \recipeServings{4--5} or \recipeServings{6 (one burrito each)}
% Place directly after the first \makeRecipeTitle on page 1.
\newcommand{\recipeServings}[1]{
    {\noindent \itshape Serves #1 \par}
    \vspace{-0.75em}
}
% --- PREP ICONS ---
% Usage: \item[\veg] Red Bell Pepper – 1, sliced
%        \item[\measure] Fish Sauce – 1 1/2 tbsp (22ml)
% Replace the bullet on lines to handle before cooking starts: produce to cut
% (carrot), and spices, seasonings, sauces, and pastes with a measured amount
% (measuring spoon, drawn with TikZ). "To taste" lines get neither.
\usepackage{fontawesome5}
\usepackage{tikz}
\newcommand{\veg}{\raisebox{0.25ex}{\small\faCarrot}}
\newcommand{\measure}{\tikz[x=0.9em, y=0.9em,
        baseline={([yshift=-0.36em]current bounding box.center)}]{
    \begin{scope}[rotate=45]
        % Bowl with a rim highlight, and a flat handle with a hanging hole.
        \fill (0,0) ellipse [x radius=0.33, y radius=0.29];
        \draw[white, line width=0.035em, line cap=round]
            (-0.1,0.173) arc[start angle=120, end angle=210, radius=0.2];
        \fill (0.2,-0.07) -- (0.86,-0.105) arc[start angle=-90, end angle=90, radius=0.105]
            -- (0.2,0.07) -- cycle;
        \fill[white] (0.86,0) circle (0.04);
    \end{scope}
}}
% --- CUSTOM COMMAND FOR SOURCE CREDIT ---
% Usage: \recipeSource{Creator -- Video Title}{https://...}
% Place just before \end{document}; pins the credit to the bottom of the last page.
\usepackage[hidelinks]{hyperref}
\urlstyle{same}
\newcommand{\recipeSource}[2]{
    \vfill
    {\noindent \footnotesize \textbf{Source:} #1 \\ \url{#2}}
}


\begin{document}

% ================= PAGE 1 =================
\makeRecipeTitle{Frozen Chicken Stock}
\recipeYield{about 12 cups (3 quarts)}

\section*{Ingredients}

\begin{itemize}
      \item Chicken Carcasses – 2, plus the necks, hearts, gizzards, thigh bones, and wing tips
      \item[\veg] Yellow Onions – 2, quartered (skin on)
      \item[\veg] Carrots – 2, cut into thirds
      \item[\veg] Celery Ribs – 2, cut into thirds
      \item[\veg] Garlic Cloves – 6, smashed
      \item[\measure] Black Peppercorns – 1 tsp
      \item Bay Leaves – 2
      \item Cold Water – about 16 cups (4 quarts), enough to cover by 1 inch
\end{itemize}

\section*{Notes}

\begin{itemize}
      \item \textbf{No Salt:} Leave the stock unsalted. It gets reduced and seasoned in whatever recipe it goes into, so salting now makes those dishes too salty.
      \item \textbf{Skip the Livers:} They make stock bitter and cloudy. Everything else from the giblet bag goes in.
      \item \textbf{Not Ready Yet?} Bag the carcasses and freeze them. They go into the pot straight from frozen; just add 30 minutes to the simmer.
      \item \textbf{Pot Size:} You need at least an 8-quart stockpot. If the bones don't fit, break the carcasses in half with your hands.
      \item \textbf{Storage:} Keeps 4 days in the fridge or 6 months in the freezer.
\end{itemize}

\newpage
\makeRecipeTitle{Frozen Chicken Stock}
\section*{Instructions}

\begin{enumerate}

      \item \textbf{Load the Pot} \\
            Put the carcasses and bones in the stockpot and add cold water until they're covered by 1 inch. Starting cold pulls more flavor from the bones and makes the foam easier to skim.

      \item \textbf{Bring to a Simmer} \\
            Heat over medium-high until small bubbles just start breaking the surface, then turn it down to low. For the first 20 minutes, skim off the gray foam that rises. \textbf{Don't let it boil:} boiling turns the stock cloudy and greasy.

      \item \textbf{Add the Aromatics} \\
            Add the onions, carrots, celery, garlic, peppercorns, and bay leaves. Simmer uncovered at the lowest bubble your stove holds (3--4 hours). If the bones start poking out, add a little hot water.

      \item \textbf{Strain} \\
            Ladle the stock through a fine-mesh strainer into a large bowl and throw out the solids. The bones have given up everything by now.

      \item \textbf{Cool and Skim} \\
            Set the bowl in a sink of ice water and stir until the stock is just warm (about 20 minutes), then refrigerate overnight. In the morning, lift off the hardened fat cap. The stock underneath should jiggle like jelly, which means it's full of body.

      \item \textbf{Portion and Freeze} \\
            Ladle into 1-cup and 4-cup containers or zip-top bags and label each one with the amount and date. Freeze bags flat on a sheet pan so they stack.

\end{enumerate}

\end{document}
